\documentclass[a4paper]{article}

% Basic packages
\usepackage[fontset=fandol]{ctex}
\usepackage{amsmath, amssymb}
\usepackage{graphicx}
\usepackage[margin=2.5cm]{geometry}
\usepackage[most]{tcolorbox}
\usepackage{etoolbox}
\usepackage{listings}
\usepackage{hyperref}
\usepackage{booktabs}
\usepackage{subcaption}
\usepackage{float}
\usepackage{tikz}
\usetikzlibrary{positioning}
\IfFileExists{pgfplots.sty}{
    \usepackage{pgfplots}
    \pgfplotsset{compat=1.18}
}{}

% Highlight boxes
\newtcolorbox{knowledgebox}[1]{
    enhanced, colback=blue!5!white, colframe=blue!75!black, colbacktitle=blue!75!black,
    coltitle=white, fonttitle=\bfseries, title=#1, attach boxed title to top left={yshift=-2mm, xshift=2mm},
    boxrule=1pt, sharp corners
}

\newtcolorbox{importantbox}[1]{
    enhanced, colback=yellow!10!white, colframe=yellow!80!black, colbacktitle=yellow!80!black,
    coltitle=black, fonttitle=\bfseries, title=#1, sharp corners
}

\newtcolorbox{warningbox}[1]{
    enhanced, colback=red!5!white, colframe=red!75!black, colbacktitle=red!75!black,
    coltitle=white, fonttitle=\bfseries, title=#1, sharp corners
}

% Listings
\lstset{
    language=Python,
    basicstyle=\ttfamily\small,
    keywordstyle=\color{blue},
    stringstyle=\color{red!60!black},
    commentstyle=\color{green!60!black},
    breaklines=true,
    frame=single,
    numbers=left,
    numberstyle=\tiny\color{gray},
    captionpos=b,
    extendedchars=false
}

% Front-page metadata
\newcommand{\notetitle}{CS224N Lecture 12: Efficient Neural Network Training}
\newcommand{\noteauthors}{整理：AI Course Notes Contributors；授课：Chris Manning}
\newcommand{\notedate}{2024年5月}
\newcommand{\videochannel}{Stanford Online}
\newcommand{\videopublishdate}{2024}
\newcommand{\videoduration}{约55分钟}
\newcommand{\videourl}{https://www.youtube.com/watch?v=UVX7SYGCKkA}
\newcommand{\videocoverpath}{}
\newcommand{\repourl}{https://github.com/hqhq1025/ai-course-notes}


% Footer with repo info
\usepackage{fancyhdr}
\pagestyle{fancy}
\fancyhf{}
\fancyfoot[C]{\thepage}
\fancyfoot[R]{\footnotesize\color{gray}\href{\repourl}{AI Course Notes}}
\renewcommand{\headrulewidth}{0pt}
\renewcommand{\footrulewidth}{0pt}

\begin{document}
\setlength{\emergencystretch}{2em}

\begin{titlepage}
\centering
{\Large 课程笔记\par}
\vspace{1.2cm}
{\huge\bfseries \notetitle\par}
\vspace{0.8cm}
{\large \noteauthors\par}
\vspace{0.3cm}
{\large \notedate\par}
\vspace{1.2cm}

\vspace{1.2cm}
\vfill
\begin{tcolorbox}[width=0.9\textwidth, colback=blue!3!white, colframe=blue!55!black, sharp corners]
\textbf{课程版本：} 2024 年中文讲义整理版\par
\textbf{笔记来源：} AI Course Notes Contributors\par
\textbf{本版处理：} 移除课件截图和对应图注，不包含 Stanford 原始课件、图片或视频。\par
\textbf{授权：} 笔记依 CC BY-NC-SA 4.0 分享；完整许可与改动记录见下一页。
\end{tcolorbox}
\end{titlepage}

\begin{center}
\fbox{\begin{minipage}{0.91\textwidth}
\small
\textbf{来源与许可}\\
本中文笔记由 AI Course Notes Contributors 整理，课程版本为 2024；源稿：\url{https://github.com/hqhq1025/ai-course-notes/tree/717e2df65c3d2eee593c171156a30d4f8f6812b9/cs224n/lecture12}。\\
笔记内容依 CC BY-NC-SA 4.0 分享：\url{https://creativecommons.org/licenses/by-nc-sa/4.0/}。\\
本副本的改动：移除 Stanford 原始课件截图与依赖这些截图的图注，移除课程视频封面/链接，并清理 TeX 源稿中的行尾空白后重新编译为可独立阅读的 PDF；同时加入本来源与许可说明。完整许可文本随本文件夹提供。Stanford 课件、课件图片和视频不包含在该授权内；请从对应年份 Stanford 课程页访问原始课件。
\end{minipage}}
\end{center}
\clearpage

\tableofcontents
\newpage

%% ========== Part 1: 混合精度训练 ==========

\section{引言：高效训练神经网络的重要性}

本节课的主题是高效神经网络训练（Efficient Neural Network Training），内容涵盖三大主题：（1）混合精度训练，（2）多GPU训练（DDP与FSDP），（3）参数高效微调（LoRA）。这些技术对于实际项目中的模型训练至关重要。


% Raster figure omitted from this study copy.



\begin{importantbox}{本课三大核心主题}
\begin{enumerate}
    \item \textbf{Mixed Precision Training}（混合精度训练）：用更少的比特表示参数和梯度，节省内存并加速计算
    \item \textbf{Multi-GPU Training}（多GPU训练）：DDP 和 FSDP（ZeRO Stage 1--3）的原理与实践
    \item \textbf{Parameter Efficient Finetuning}（参数高效微调）：LoRA 的核心思想与实现
\end{enumerate}
\end{importantbox}

到目前为止，CS224N 课程已经讲解了如何将词转化为向量、如何将句子转化为向量、Transformer 架构、预训练等内容。本节课转向一个不同但同样重要的方向——如何在 GPU 上高效地训练大型模型。虽然这些内容``与自然语言没有直接关系''，但对于最终项目（和实际工作）极其有用。

\subsection{本章小结}

高效训练技术是将理论模型落地到实际GPU硬件上的桥梁。混合精度训练节省内存和时间，多GPU训练突破单卡容量限制，参数高效微调让资源受限的研究者也能适配大模型。

%% ========== 浮点数表示 ==========

\section{浮点数表示基础}

\subsection{FP32：标准单精度浮点数}

在深度学习中，模型参数和梯度通常以浮点数的形式存储在GPU显存中。FP32（32位浮点数）是最基本的数据类型，每个参数占用 \textbf{4字节}。


% Raster figure omitted from this study copy.



FP32 的数值由以下公式计算：

$$(-1)^{B} \times 2^{E-127} \times \left(1 + \sum_{i=1}^{23} b_{23-i} \cdot 2^{-i}\right)$$

\begin{itemize}
    \item $B$：符号位（sign），决定正负
    \item $E$：指数位（exponent，8位），决定\textbf{动态范围}——能表示多大或多小的数
    \item $b_i$：尾数位（mantissa，23位），决定\textbf{精度}——数值的精确程度
\end{itemize}


% Raster figure omitted from this study copy.



\begin{knowledgebox}{指数位与尾数位的权衡}
\textbf{指数位越多}，能表示的数的范围越大（更小的数和更大的数都能表示）。\\
\textbf{尾数位越多}，数值的精度越高（相邻可表示数之间的间隔更小）。\\
在总位数固定的前提下，指数位与尾数位之间存在权衡——这是理解FP16、BF16等格式的关键。
\end{knowledgebox}

\subsection{FP16：半精度浮点数}

FP16 只有16位（2字节），内存需求是 FP32 的一半。其格式为：1位符号位 + 5位指数位 + 10位尾数位。


% Raster figure omitted from this study copy.



\begin{table}[H]
\centering
\begin{tabular}{lccc}
\toprule
\textbf{数据类型} & \textbf{总位数} & \textbf{指数位} & \textbf{尾数位} \\
\midrule
FP32  & 32 & 8  & 23 \\
FP16  & 16 & 5  & 10 \\
BF16  & 16 & 8  & 7  \\
\bottomrule
\end{tabular}
\caption{三种常见浮点数格式的位数分配}
\end{table}

FP16 的两个核心问题：

\begin{enumerate}
    \item \textbf{动态范围不足}：指数位仅5位，导致非常小的数会被\textbf{下溢为零}（underflow）。例如，小于约 $6 \times 10^{-5}$ 的数在 FP16 中直接变为0。
    \item \textbf{精度不足}：尾数位仅10位，导致\textbf{舍入误差}。例如，1.0001 在 FP16 中会被舍入为 1.0。
\end{enumerate}


% Raster figure omitted from this study copy.

\footnotetext{Slides 第11页。图来自 NVIDIA 博客。}

\begin{warningbox}{FP16 梯度下溢是致命问题}
在神经网络训练中，梯度通常非常小。如果直接用 FP16 计算梯度，\textbf{超过一半的梯度值会被四舍五入为零}，导致模型无法有效学习。这就是为什么不能简单地将所有计算切换到 FP16。
\end{warningbox}

\subsection{本章小结}

浮点数的位数分配（指数位 vs 尾数位）决定了动态范围和精度之间的权衡。FP32 是标准格式（4字节/参数），FP16 节省一半内存但存在严重的动态范围和精度问题。理解这些限制是混合精度训练的前提。

%% ========== 混合精度训练 ==========

\section{混合精度训练}

\subsection{初步方案：FP32 主权重 + FP16 计算}

既然纯 FP16 有问题，一个自然的想法是：\textbf{保留一份 FP32 的模型副本（Master Weights），但在前向和反向传播中使用 FP16 计算}。


% Raster figure omitted from this study copy.



具体步骤：
\begin{enumerate}
    \item 维护 FP32 格式的模型参数副本（Master Weights）
    \item 将 Master Weights 转换为 FP16，执行前向传播
    \item 在 FP16 中计算梯度（反向传播）
    \item 将 FP16 梯度上转为 FP32
    \item 用 FP32 梯度更新 Master Weights
    \item 将更新后的 Master Weights 复制回 FP16 版本
\end{enumerate}

\begin{warningbox}{梯度下溢问题仍未解决}
虽然权重更新的精度问题通过 FP32 Master Weights 解决了，但反向传播中计算的\textbf{梯度本身仍然是 FP16 的}。那些极小的梯度值仍然会在 FP16 中下溢为零，丢失重要的学习信号。
\end{warningbox}

\subsection{Loss Scaling：解决梯度下溢}

为了避免 FP16 下的梯度下溢，引入 \textbf{Loss Scaling} 技巧：在计算梯度之前，将损失函数乘以一个较大的常数（如1000），使得梯度也被等比放大。放大后的梯度值更大，不容易下溢为零。在 FP32 中更新权重之前，再将梯度除以该常数。


% Raster figure omitted from this study copy.



\begin{importantbox}{混合精度训练完整方案（FP16 + Loss Scaling）}
\begin{enumerate}
    \item 维护 FP32 格式的 Master Weights
    \item 在 FP16 中执行前向传播
    \item \textbf{将损失乘以一个大常数}（Loss Scaling），人为放大梯度
    \item 在 FP16 中计算梯度（此时梯度已被放大，不容易下溢）
    \item 将梯度转为 FP32 后\textbf{除以 Scaling Factor}
    \item 用 FP32 更新 Master Weights
    \item 复制回 FP16 版本
\end{enumerate}
\end{importantbox}

在 PyTorch 中，使用 \texttt{torch.cuda.amp.GradScaler} 和 \texttt{torch.cuda.amp.autocast} 可以方便地实现此方案：

\begin{lstlisting}[caption={PyTorch 中的 FP16 混合精度训练（含 GradScaler）}]
scaler = torch.cuda.amp.GradScaler()

for input, target in data:
    optimizer.zero_grad()
    with torch.cuda.amp.autocast():
        output = model(input)
        loss = loss_fn(output, target)
    scaler.scale(loss).backward()
    scaler.step(optimizer)
    scaler.update()
\end{lstlisting}

\begin{warningbox}{Loss Scaling 的局限性}
Loss Scaling 的常数需要精心选择：太小则梯度仍然下溢，太大则梯度上溢为 NaN。\texttt{GradScaler} 通过动态调整 Scaling Factor 来应对不同训练阶段的梯度分布变化，但这增加了训练的复杂性。
\end{warningbox}

\subsection{BFloat16：更优雅的解决方案}

Loss Scaling 的根本原因是 FP16 的动态范围不足（仅5位指数）。\textbf{BFloat16}（Brain Float 16，由 Google Brain 提出）提供了一个更根本的解决方案：


% Raster figure omitted from this study copy.



BF16 的设计理念：

\begin{itemize}
    \item \textbf{8位指数}：与 FP32 \textbf{完全相同}的动态范围，彻底消除下溢问题
    \item \textbf{7位尾数}：精度比 FP16（10位）更低，但实验表明对神经网络训练影响极小
    \item \textbf{无需 Loss Scaling}：因为动态范围足够，不再需要人为放大梯度
\end{itemize}


% Raster figure omitted from this study copy.



使用 BF16 的代码极其简洁：

\begin{lstlisting}[caption={PyTorch 中的 BF16 混合精度训练（无需 GradScaler）}]
# Creates model and optimizer in default precision
model = Net().cuda()
optimizer = optim.SGD(model.parameters(), ...)

for input, target in data:
    optimizer.zero_grad()
    # Enables autocasting for the forward pass
    with torch.autocast(device_type="cuda"):
        output = model(input)
        loss = loss_fn(output, target)
    # No GradScaler needed!
    loss.backward()
    optimizer.step()
\end{lstlisting}

\begin{importantbox}{BF16 是当前最佳实践}
如果你的 GPU 支持 BF16（Ampere 架构及更新：A100、H100、A6000 等），\textbf{始终使用 BF16 而非 FP16}。BF16 无需 Loss Scaling，代码更简洁，训练更稳定。可通过 \texttt{torch.cuda.is\_bf16\_supported()} 检查支持情况。
\end{importantbox}

\subsection{混合精度训练的实际效果}

以在单块 A100 上微调 DistilBERT 进行情感分类为例：

\begin{table}[H]
\centering
\begin{tabular}{lccc}
\toprule
\textbf{数据类型} & \textbf{训练时间} & \textbf{准确率} & \textbf{显存占用} \\
\midrule
Float64      & $\sim$25 min & 高 & 最大 \\
Float32      & 基准     & 基准 & 基准 \\
Mixed (BF16) & \textbf{减少约1/3} & \textbf{略有提升} & \textbf{显著减少} \\
\bottomrule
\end{tabular}
\caption{不同精度下 DistilBERT 微调的性能对比}
\end{table}

BF16 混合精度训练不仅节省了约 1/3 的训练时间和大量显存，准确率甚至略有提升——这是因为低精度表示具有一定的\textbf{正则化效果}。训练加速的根本原因是：半精度矩阵乘法在现代 GPU 的 Tensor Core 上速度更快。

\subsection{本章小结}

\begin{itemize}
    \item FP16 节省内存但存在梯度下溢和精度问题，需要 Loss Scaling
    \item BF16 保留 FP32 的动态范围，牺牲精度，无需 Loss Scaling
    \item \textbf{混合精度训练应始终启用}：几乎不影响模型质量，显著节省时间和内存
    \item 新架构 GPU 优先使用 BF16，旧架构使用 FP16 + GradScaler
\end{itemize}

%% ========== Part 2: 多GPU训练 ==========

\section{多GPU训练：DDP 与 ZeRO}

\subsection{GPU 显存占用分析}

在讨论多GPU训练之前，需要理解单GPU上的显存都用在了哪里。以混合精度训练为例，每个参数的显存开销为：


% Raster figure omitted from this study copy.



\begin{table}[H]
\centering
\begin{tabular}{llc}
\toprule
\textbf{组件} & \textbf{精度} & \textbf{每参数字节数} \\
\midrule
模型参数     & FP16 & 2 \\
梯度        & FP16 & 2 \\
Master Weights（FP32 副本） & FP32 & 4 \\
Adam 动量（momentum）       & FP32 & 4 \\
Adam 方差（variance）       & FP32 & 4 \\
\midrule
\textbf{合计} & & \textbf{16} \\
\bottomrule
\end{tabular}
\caption{混合精度 + Adam 优化器的每参数显存开销}
\end{table}

\begin{knowledgebox}{优化器状态的显存开销}
很多人第一次了解到优化器也需要大量显存时会感到惊讶。使用 Adam 优化器时，每个参数需要额外存储动量（momentum）和方差（variance），各占 4 字节（FP32）。加上 FP32 Master Weights，仅优化器状态就需要 \textbf{12字节/参数}——是模型参数本身（2字节）的 6 倍！
\end{knowledgebox}

\subsection{分布式数据并行（DDP）}

\textbf{Distributed Data Parallel (DDP)} 是最基本的多GPU训练策略。


% Raster figure omitted from this study copy.



DDP 的工作流程：
\begin{enumerate}
    \item 将数据集均匀分割到 $N$ 个 GPU 上（每个 GPU 处理 $1/N$ 的数据）
    \item 每个 GPU 维护\textbf{完整的}模型副本和优化器状态
    \item 每个 GPU 独立执行前向传播和反向传播，得到各自的梯度
    \item 通过 \textbf{All-Reduce} 操作同步梯度：汇总所有 GPU 的梯度并分发给所有 GPU
    \item 每个 GPU 用汇总后的梯度更新自己的模型，保持同步
\end{enumerate}


% Raster figure omitted from this study copy.



\begin{importantbox}{All-Reduce：DDP 的核心原语}
\textbf{All-Reduce} 是一个 MPI（消息传递接口）原语，它将所有 GPU 上的数据进行归约（如求和），然后将结果广播到每个 GPU。在 DDP 中，All-Reduce 用于汇总梯度。通信开销为每参数 2 字节（FP16 梯度）。
\end{importantbox}

\subsection{DDP 的内存问题}

DDP 虽然简单有效，但存在严重的\textbf{内存冗余}问题：


% Raster figure omitted from this study copy.



以一个 7.5B 参数的模型（$\Psi = 7.5\text{B}$）为例，使用 Adam + 混合精度：

$$\text{每GPU显存} = (2 + 2 + K) \times \Psi = (2 + 2 + 12) \times 7.5\text{B} = 120 \text{GB}$$

其中 $K = 12$（FP32 Master Weights 4B + Adam momentum 4B + Adam variance 4B）。A100 有 80GB 显存，\textbf{连一块卡都放不下}。而且，每个 GPU 上的优化器状态完全相同——这是巨大的浪费。

\begin{warningbox}{DDP 无法解决大模型问题}
DDP 不减少任何单GPU的显存占用——它只是通过数据并行来加速训练。如果单个GPU放不下模型+优化器状态，DDP 无能为力。我们需要更聪明的策略来\textbf{分摊}显存开销。
\end{warningbox}

\subsection{ZeRO Stage 1：分片优化器状态}

\textbf{ZeRO}（Zero Redundancy Optimizer）由微软 DeepSpeed 项目提出，核心思想是将冗余的状态\textbf{分片（shard）}存储到不同 GPU 上。

ZeRO Stage 1 只分片\textbf{优化器状态}（绿色部分）：


% Raster figure omitted from this study copy.



Stage 1 的工作流程：
\begin{enumerate}
    \item 每个 GPU 持有完整的模型参数（FP16）和完整的梯度
    \item 但只持有 $1/N$ 的优化器状态
    \item 反向传播后，通过 \textbf{Reduce-Scatter} 操作，每个 GPU 只接收自己负责的参数分片的汇总梯度
    \item 每个 GPU 用本地优化器状态更新自己负责的参数分片
    \item 通过 \textbf{All-Gather} 操作，收集所有 GPU 更新后的参数，恢复完整模型
\end{enumerate}

\begin{importantbox}{ZeRO Stage 1 是免费的午餐}
关键洞察：\textbf{All-Reduce = Reduce-Scatter + All-Gather}。DDP 需要执行一次 All-Reduce，而 ZeRO Stage 1 执行 Reduce-Scatter + All-Gather，通信量完全相同。因此，ZeRO Stage 1 在\textbf{不增加任何通信开销}的情况下节省了优化器状态的显存——你应该始终使用它。
\end{importantbox}

\begin{knowledgebox}{三个 MPI 原语}
\begin{itemize}
    \item \textbf{All-Reduce}：所有 GPU 的数据归约后广播到所有 GPU（用于 DDP）
    \item \textbf{Reduce-Scatter}：所有 GPU 的数据归约后，结果的不同分片散发到不同 GPU
    \item \textbf{All-Gather}：收集所有 GPU 的分片，拼接成完整数据广播到所有 GPU
\end{itemize}
核心恒等式：All-Reduce $\equiv$ Reduce-Scatter + All-Gather
\end{knowledgebox}

\subsection{ZeRO Stage 2：分片梯度}

Stage 2 在 Stage 1 的基础上，进一步分片\textbf{梯度}。

核心技巧：\textbf{永远不实例化完整的梯度向量}。反向传播是逐层进行的——当计算完第 $j$ 层的梯度后，立即通过 Reduce 操作将该层梯度发送给负责该层的 GPU，然后释放本地的临时梯度内存。


% Raster figure omitted from this study copy.



Stage 2 的通信量也与 DDP 相当（Reduce + All-Gather $\approx$ All-Reduce），因此也是\textbf{基本免费的}显存优化。

\begin{importantbox}{ZeRO Stage 1 和 2 应始终启用}
ZeRO Stage 1（分片优化器状态）和 Stage 2（分片梯度）都不增加显著的通信开销，却能大幅减少每GPU的显存占用。在多GPU训练中，\textbf{没有理由不使用它们}。
\end{importantbox}

\subsection{ZeRO Stage 3 / FSDP：分片模型参数}

当模型大到连参数本身都放不进单个 GPU 时，就需要 \textbf{ZeRO Stage 3}，也称为 \textbf{FSDP（Fully Sharded Data Parallel）}。


% Raster figure omitted from this study copy.



Stage 3 连模型参数也进行分片：
\begin{enumerate}
    \item 将模型划分为多个 \textbf{FSDP 单元}
    \item 将每个单元的参数展平（Flat Parameter），分配到不同 GPU
    \item \textbf{前向传播时}：在计算某层之前，通过 All-Gather 临时收集该层的完整参数；计算完成后丢弃非本地分片
    \item \textbf{反向传播时}：再次 All-Gather 获取完整参数，计算梯度后通过 Reduce-Scatter 分发梯度
    \item 每个 GPU 用本地优化器状态更新本地参数分片
\end{enumerate}


% Raster figure omitted from this study copy.



\begin{warningbox}{Stage 3 有额外的通信开销}
与 Stage 1/2 不同，Stage 3 需要在前向传播和反向传播中都执行 All-Gather（分别收集参数），加上反向传播中的 Reduce-Scatter。总通信量为 \textbf{2次 All-Gather + 1次 Reduce-Scatter}，显著高于 DDP。但如果模型放不进单个 GPU，这是唯一的选择。
\end{warningbox}


% Raster figure omitted from this study copy.



好消息是：当模型足够大时，前向传播一层的计算时间足以\textbf{预取}（prefetch）下一层的参数。PyTorch 的 FSDP 默认启用这种计算-通信重叠，使得通信开销在实际中被大幅隐藏。

\subsection{GPU 显存中被忽视的部分：模型激活值}

在前面的分析中，我们忽略了一个重要的显存占用来源：\textbf{模型激活值}（Activations）。


% Raster figure omitted from this study copy.



反向传播需要保存前向传播中每一层的激活值，其显存占用与 \textbf{batch size 线性相关}。这解释了为什么增大 batch size 时容易遇到 OOM 错误。ZeRO 的三个阶段都\textbf{不处理}激活值分片——要解决激活值的显存问题，需要使用 \textbf{Activation Checkpointing}（梯度检查点）技术。

\begin{knowledgebox}{Activation Checkpointing}
Activation Checkpointing（也叫 Gradient Checkpointing）的核心思想：不保存所有层的激活值，只保存部分``检查点''层的激活值。反向传播到某层时，从最近的检查点重新计算中间激活值。这是用\textbf{额外的计算时间}换取\textbf{更少的显存占用}。
\end{knowledgebox}

\subsection{本章小结}

\begin{itemize}
    \item \textbf{DDP}：数据并行，每个 GPU 维护完整的模型/梯度/优化器状态，通过 All-Reduce 同步梯度
    \item \textbf{ZeRO Stage 1}：分片优化器状态，通信量与 DDP 相同（免费优化）
    \item \textbf{ZeRO Stage 2}：分片优化器状态 + 梯度，通信量与 DDP 相当（基本免费）
    \item \textbf{ZeRO Stage 3 / FSDP}：分片一切（参数+梯度+优化器状态），有额外通信开销，但支持超大模型
    \item 模型激活值的显存与 batch size 线性相关，需要 Activation Checkpointing 另行处理
\end{itemize}

%% ========== Part 3: 参数高效微调 ==========

\section{参数高效微调：LoRA}

\subsection{为什么需要参数高效微调}

当以下所有手段都无法解决 OOM 问题时——混合精度训练、ZeRO Stage 3、Activation Checkpointing、batch size 已经为 1——就需要考虑\textbf{参数高效微调}（Parameter-Efficient Fine-Tuning，PEFT）。


% Raster figure omitted from this study copy.



参数高效微调的动机不仅仅是显存限制：


% Raster figure omitted from this study copy.



\begin{knowledgebox}{参数高效微调的三大动机}
\begin{enumerate}
    \item \textbf{显存限制}：GPT-3 有 1750 亿参数，全量微调需要海量显存
    \item \textbf{环境与成本}：训练 GPT-3 的碳排放相当于运行燃煤电厂10小时；AI计算需求即将超过全球计算产能
    \item \textbf{科学动机}：大模型严重过参数化，对于小数据集的下游任务，参数高效微调可能泛化更好（起到正则化作用）
\end{enumerate}
\end{knowledgebox}

\subsection{全量微调 vs 参数高效微调}

\textbf{全量微调}（Full Fine-Tuning）：更新模型的所有参数 $\theta$，寻找最优的 $\Delta\theta$。

$$\theta_{\text{new}} = \theta_{\text{pretrained}} + \Delta\theta, \quad |\Delta\theta| = |\theta|$$

\textbf{参数高效微调}：只更新一小部分参数，搜索空间远小于全量微调。

$$\theta_{\text{new}} = \theta_{\text{pretrained}} + \Delta\theta, \quad |\Delta\theta| \ll |\theta|$$


% Raster figure omitted from this study copy.



参数高效微调的额外好处：每个任务只需存储一个很小的 $\Delta\theta$（而非完整的模型副本），在多任务场景下极其高效。

\subsection{LoRA：低秩适配}

\textbf{LoRA}（Low-Rank Adaptation）是当前最流行的参数高效微调方法。它基于一个经验观察：\textbf{大模型微调时的梯度具有低内在秩}（low intrinsic rank）。


% Raster figure omitted from this study copy.



对于一个预训练权重矩阵 $W \in \mathbb{R}^{d \times k}$，LoRA 将更新约束为低秩形式：

$$W_{\text{new}} = W + \alpha \cdot B \cdot A$$

\begin{itemize}
    \item $W \in \mathbb{R}^{d \times k}$：预训练权重（\textbf{冻结}，不参与训练）
    \item $A \in \mathbb{R}^{r \times k}$：LoRA 矩阵 A（可训练）
    \item $B \in \mathbb{R}^{d \times r}$：LoRA 矩阵 B（可训练）
    \item $r$：秩，远小于 $d$ 和 $k$（典型值：$r = 8$）
    \item $\alpha$：缩放系数，控制预训练知识与新任务知识的权衡
\end{itemize}

\begin{importantbox}{LoRA 的参数节省}
原始权重矩阵 $W$ 有 $d \times k$ 个参数。LoRA 只需学习 $A$（$r \times k$ 个参数）和 $B$（$d \times r$ 个参数），共 $r \times (d + k)$ 个参数。当 $r \ll \min(d, k)$ 时，可训练参数量远小于原始参数量。例如，$d = k = 768$，$r = 8$ 时，可训练参数量仅为原来的 $\frac{2 \times 8}{768} \approx 2\%$。
\end{importantbox}

\subsection{LoRA 的实现}

在代码层面，LoRA 的实现非常简洁：


% Raster figure omitted from this study copy.



\begin{lstlisting}[caption={LoRA 的 PyTorch 实现}]
input_dim = 768   # hidden size of pre-trained model
output_dim = 768  # output size of the layer
rank = 8          # rank 'r' for low-rank adaptation

W = ...  # from pretrained network, shape: input_dim x output_dim

W_A = nn.Parameter(torch.empty(input_dim, rank))  # LoRA weight A
W_B = nn.Parameter(torch.empty(rank, output_dim)) # LoRA weight B

# Initialization
nn.init.kaiming_uniform_(W_A, a=math.sqrt(5))
nn.init.zeros_(W_B)  # B initialized to zero!

def regular_forward_matmul(x, W):
    h = x @ W
    return h

def lora_forward_matmul(x, W, W_A, W_B):
    h = x @ W           # regular matrix multiplication (frozen)
    h += x @ (W_A @ W_B) * alpha  # add scaled LoRA weights
    return h
\end{lstlisting}

\begin{knowledgebox}{LoRA 初始化策略}
$B$ 矩阵初始化为\textbf{全零}，$A$ 矩阵使用 Kaiming 初始化。这意味着训练开始时 $\Delta W = B \cdot A = 0$，模型行为与预训练模型完全一致。随着训练进行，$B$ 逐渐偏离零，模型开始适配下游任务。这种初始化确保了训练的稳定性。
\end{knowledgebox}

\subsection{LoRA 的关键超参数}


% Raster figure omitted from this study copy.



\textbf{1. 应用位置}：LoRA 应应用于 Self-Attention 中的\textbf{Query 矩阵（$W_q$）和 Value 矩阵（$W_v$）}，这是经验上效果最好的选择。


% Raster figure omitted from this study copy.



\textbf{2. 秩 $r$}：实验表明，即使 $r$ 非常小（$r = 1$ 或 $r = 4$），LoRA 也能获得很好的性能。$r = 8$ 是一个推荐的起始值。

\textbf{3. 缩放系数 $\alpha$}：通常设为 1。$\alpha < 1$ 表示更偏向保留预训练知识，$\alpha > 1$ 表示更偏向学习新任务知识。

\begin{importantbox}{LoRA 实用指南}
\begin{itemize}
    \item 应用于 \textbf{Query 和 Value 矩阵}
    \item 秩设为 \textbf{$r = 8$}
    \item $\alpha$ 设为 \textbf{1}
    \item 这些默认值在大多数场景下效果很好，可以作为起点
\end{itemize}
\end{importantbox}

\subsection{LoRA 的优势}


% Raster figure omitted from this study copy.



\begin{enumerate}
    \item \textbf{显存高效}：只需存储和优化两个小矩阵 $A$ 和 $B$，大幅减少梯度和优化器状态的显存
    \item \textbf{推理无开销}：训练完成后，可以将 $\alpha \cdot B \cdot A$ 直接合并到 $W$ 中，推理时与原模型无异
    \item \textbf{多任务高效}：每个任务只需存储一组小的 $A, B$ 矩阵；切换任务时只需替换 LoRA 权重
    \item \textbf{收敛于全量微调}：当 $r$ 增大到等于 $\min(d, k)$ 时，LoRA 等价于全量微调——提供了一个连续的控制旋钮
    \item \textbf{正则化效果}：低秩约束天然地限制了搜索空间，在小数据集上可能泛化更好
\end{enumerate}

\subsection{本章小结}

\begin{itemize}
    \item LoRA 将权重更新约束为低秩矩阵乘积 $\Delta W = B \cdot A$，大幅减少可训练参数
    \item 应用于 Self-Attention 的 $W_q$ 和 $W_v$ 矩阵，秩 $r = 8$，$\alpha = 1$ 是好的起点
    \item 推理时无额外开销，多任务时只需存储小矩阵
    \item 在资源受限环境下，LoRA 是微调大模型的首选方案
\end{itemize}

%% ========== 实用决策流程 ==========

\section{实用决策流程图}

Shikhar Murty 在课程结尾给出了一个非常实用的决策流程图，帮助你在实际项目中选择合适的训练策略：


% Raster figure omitted from this study copy.

\footnotetext{Slides 第65页（最后一页）。}

\begin{importantbox}{高效训练决策流程}
\begin{enumerate}
    \item \textbf{始终使用混合精度训练}。如果 GPU 支持 Ampere 架构（A100/H100/A6000），使用 BF16；否则使用 FP16 + GradScaler
    \item 尝试 \textbf{batch size = 1}。如果能放下：
    \begin{itemize}
        \item 增大 batch size
        \item 始终使用 \textbf{ZeRO Stage 2}（免费优化）
    \end{itemize}
    \item 如果 batch size = 1 也放不下：
    \begin{itemize}
        \item 尝试 \textbf{ZeRO Stage 3 / FSDP}
        \item 尝试 \textbf{Activation Checkpointing}
    \end{itemize}
    \item 如果以上全部失败：使用 \textbf{LoRA}
    \begin{itemize}
        \item 应用于 $W_q$ 和 $W_v$
        \item 秩 $r = 8$，$\alpha = 1$
    \end{itemize}
\end{enumerate}
\end{importantbox}

\begin{warningbox}{此流程假设多 GPU 环境}
上述流程中的 ZeRO / FSDP 假设你有多个 GPU。如果只有单个 GPU，ZeRO 无法帮助你——此时应直接考虑量化（quantization）或 LoRA 等技术。
\end{warningbox}

\subsection{本章小结}

决策流程的核心逻辑是``从简单到复杂''：先用零成本的优化（混合精度、ZeRO Stage 1/2），再考虑有代价的优化（ZeRO Stage 3），最后采用改变训练范式的方法（LoRA）。在每一步都优先选择开销最小的方案。

%% ========== 总结与延伸 ==========

\section{总结与延伸}

\subsection{全课知识图谱}

本课建立了一条从数值表示到训练策略的完整认知链：

\begin{center}
\begin{tikzpicture}[node distance=0.9cm, every node/.style={draw, rounded corners, minimum width=4cm, minimum height=0.8cm, font=\small}]
\node (fp) {浮点数表示};
\node (mp) [below=of fp] {混合精度训练};
\node (ddp) [below=of mp] {DDP 数据并行};
\node (zero) [below=of ddp] {ZeRO / FSDP};
\node (lora) [below=of zero] {LoRA 参数高效微调};
\draw[->, thick] (fp) -- (mp);
\draw[->, thick] (mp) -- (ddp);
\draw[->, thick] (ddp) -- (zero);
\draw[->, thick] (zero) -- (lora);
\node[draw=none, right=1cm of fp, text width=5.5cm, font=\scriptsize] {FP32/FP16/BF16 的指数位与尾数位权衡};
\node[draw=none, right=1cm of mp, text width=5.5cm, font=\scriptsize] {Master Weights + Loss Scaling / BF16 AutoCast};
\node[draw=none, right=1cm of ddp, text width=5.5cm, font=\scriptsize] {All-Reduce 同步梯度，内存无优化};
\node[draw=none, right=1cm of zero, text width=5.5cm, font=\scriptsize] {分片优化器/梯度/参数，Reduce-Scatter + All-Gather};
\node[draw=none, right=1cm of lora, text width=5.5cm, font=\scriptsize] {低秩分解 $\Delta W = BA$，冻结原参数};
\end{tikzpicture}
\end{center}

\subsection{关键 Takeaways}

\begin{importantbox}{五条核心原则}
\begin{enumerate}
    \item \textbf{始终启用混合精度训练}：BF16（Ampere+）或 FP16 + GradScaler，几乎没有质量损失
    \item \textbf{优化器状态是显存大户}：Adam 的 momentum 和 variance 占 12 字节/参数，远超模型本身
    \item \textbf{ZeRO Stage 1/2 免费}：利用 All-Reduce = Reduce-Scatter + All-Gather 的等价性，不增加通信即可节省显存
    \item \textbf{FSDP 支持超大模型}：有额外通信开销，但可通过计算-通信重叠来隐藏
    \item \textbf{LoRA 是最后的利器}：当全量微调不可行时，低秩适配以极少的参数达到接近全量微调的效果
\end{enumerate}
\end{importantbox}

\subsection{拓展阅读}

\begin{itemize}
    \item Micikevicius et al., 2018. \textit{Mixed Precision Training}. ICLR 2018. \url{https://arxiv.org/abs/1710.03740}
    \item Rajbhandari et al., 2020. \textit{ZeRO: Memory Optimizations Toward Training Trillion Parameter Models}. SC 2020. \url{https://arxiv.org/abs/1910.02054}
    \item Hu et al., 2021. \textit{LoRA: Low-Rank Adaptation of Large Language Models}. ICLR 2022. \url{https://arxiv.org/abs/2106.09685}
    \item PyTorch FSDP 文档：\url{https://pytorch.org/docs/stable/fsdp.html}
    \item PyTorch AMP 文档：\url{https://pytorch.org/docs/stable/amp.html}
    \item Hugging Face PEFT 库：\url{https://github.com/huggingface/peft}
    \item DeepSpeed ZeRO 文档：\url{https://www.deepspeed.ai/tutorials/zero/}
\end{itemize}

\end{document}
